\documentclass[11pt]{article}

\usepackage[margin=1in]{geometry}
\usepackage{amsmath,amssymb}
\usepackage{booktabs}
\usepackage{array}
\usepackage{graphicx}
\usepackage[section]{placeins}
\usepackage[hidelinks]{hyperref}

\renewcommand{\thetable}{S\arabic{table}}
\renewcommand{\thesection}{S\arabic{section}}
\renewcommand{\thefigure}{S\arabic{figure}}

\title{Supplementary Material\\[0.3em]\large Uncertainty-Aware Bayesian Mineral Prospectivity Mapping under Spatial Separation: Heterogeneous Predictor Relationships across the Central African Copperbelt}
\author{}
\date{}

\begin{document}
\maketitle

\setlength{\parindent}{0pt}
\setlength{\parskip}{0.5em}

All values in this supplement come from the frozen V14 benchmark artefacts or from summaries derived from them without refitting. The artefacts are listed with SHA-256 checksums in \texttt{paper\_1/benchmark/SHA256SUMS}. The tables are generated by \texttt{src/paper1\_supplement\_summaries.py}, \texttt{src/audit\_lithology\_class\_assignment.py}, and \texttt{src/build\_paper1\_supplement\_tables.py}, so their numbers are not transcribed by hand. Fold~$k$ denotes the fit in which fold~$k$ is held out; Fold~1 is at the south-eastern end of the belt.

%==========================================================================
\section{Predictor reconstruction and verification}
\label{sec:s1}

The modelling table inherits labels, domains, centroids, fault distance, gravity, and the lithology class from an earlier feature build (the V5 table). No script in the repository creates that build. Only the contact-distance column was rebuilt for the V14 benchmark, and its rebuild is scripted and checksummed. Each inherited step was reconstructed from the source products and the stored values. Table~\ref{tab:s1} summarizes the confirmation status of each step.

\begin{table}[ht]
\centering
\small
\caption{Reconstruction of the predictors and response. \emph{Confirmed}: reproduced from the source product and a documented procedure. \emph{Partly reproduced}: reproduced for a stated subset only. \emph{Inferred}: consistent with the stored values but not documented at the time of preparation.}
\label{tab:s1}
\begin{tabular}{>{\raggedright\arraybackslash}p{0.17\textwidth}>{\raggedright\arraybackslash}p{0.55\textwidth}>{\raggedright\arraybackslash}p{0.18\textwidth}}
\toprule
Variable & Reconstruction & Status \\
\midrule
Occurrence response & 187 USGS records (65 deposits, 102 prospects, 20 grouped records); 184 fall inside the 1,899-cell frame and map to 139 positive cells, 138 after complete-case filtering. 54 records lie within 250\,m of a cell boundary, 113 within 500\,m, and 162 within 1\,km; the median distance to a boundary is 398\,m. & Confirmed (point-in-polygon) \\
\addlinespace
Contact distance & Minimum planar distance from each cell centroid to the 259 internal contact lines of the polygonized surficial-lithology raster, excluding the outer tract boundary. Minimum 11.13\,m, median 11,979.61\,m, maximum 70,130.56\,m, no zero values. The contact lines plus the outer boundary reproduce the full polygon linework (partition check passed). & Confirmed (scripted, checksummed) \\
\addlinespace
Lithology class & Attribute of a nearest contact line rather than the lithology of the cell (Section~\ref{sec:s2}). & Inferred \\
\addlinespace
Fault distance & Distance from the cell centroid to the nearest of 126 lines digitized from Fig.~1 of Selley et al.\ (2005). The surviving line layer reproduces 1,672 of 1,899 stored values within 1\,m; 227 differ by more than 1\,m (26 positive cells), with a maximum difference of 33,076.98\,m. The layer has no names, types, dips, or displacements. & Partly reproduced \\
\addlinespace
Bouguer gravity & Nearest 2 arc-minute grid node to the centroid, with no interpolation; centroid-to-node offsets are 0.07--2.56\,km (median 1.47\,km). The grid header names neither the product nor the datum. The product is identified as the WGM2012 complete Bouguer anomaly from companion free-air, disturbance, and topography grids that share its limits and spacing. The source raster is not preserved; only the sampled values survive. Bouguer and isostatic values correlate at $r=0.892$; only Bouguer enters the model. & Inferred \\
\addlinespace
Tectonic domain & Labels from an 11-polygon layer digitized from Fig.~14 of Daly et al.\ (2025). The digitized polygons contain small gaps and overlaps. The stored labels of 49 complete-case background cells differ from a strict centroid-in-polygon join; no positive cell differs. 27 cells have no label (one positive) and are excluded. & Partly reproduced \\
\addlinespace
Grid support & 5\,km grid clipped to the union of five USGS permissive tracts: 1,899 cells covering 21,790.36\,km$^2$, of which 253 have the full 25\,km$^2$ area (254 at a 24.999\,km$^2$ threshold). 141 stored centroids lie outside their clipped cell polygon, including 12 positive cells. & Confirmed \\
\bottomrule
\end{tabular}
\end{table}

\paragraph{Verification of the posterior samples.} For each of the four fold fits, the posterior samples were checked against the modelling table. The training labels implied by each trace match the table, there are no divergent transitions, and across all parameters the maximum rank-normalized $\hat R$ is 1.004 and the minimum bulk effective sample size is 1,219 (ArviZ 0.23.4; per fold 1,365, 1,331, 1,219, and 1,520). The recomputed out-of-fold metrics reproduce the reported values: pooled ROC-AUC 0.7155, fold values 0.8753, 0.7357, 0.6851, and 0.6031. The post-refit validation record is \texttt{figures/audit/v14\_model\_output\_validation\_post\_refit\_20261021.json} (status \texttt{PASS}).

%==========================================================================
\section{Assignment of the surficial-lithology class}
\label{sec:s2}

The column used as the surficial-lithology covariate was inherited from the earlier feature build, and no assignment rule was recorded. Four candidate rules were tested against the source layers (Table~\ref{tab:s2}, upper panel). The stored class rarely equals the lithology at the centroid or the majority lithology in the cell. It does equal, in 95.4\% of cells, the class attribute of one of the contact lines at the minimum distance from the cell polygon. That set of tied contacts is the output of a nearest-feature join from cell polygons to contact lines. Contact lines meet at shared nodes, so the nearest point is usually shared by several lines. In 1,890 of 1,899 cells, these tied contacts carry different classes, and the stored value is one of them, chosen by an undocumented rule. The 87 cells not matched by any tied contact include 5 positive cells. They may come from an earlier version of the contact layer.

The covariate is therefore an attribute of the nearest contact, chosen between the two sides of that contact, rather than the lithology of the cell. The manuscript describes the lithology terms accordingly. Replacing it with the centroid or majority class would change the model and would need a separate, versioned refit; this was not done for the reported benchmark.

The lower panel of Table~\ref{tab:s2} lists the indicator coding. Metasedimentary is the reference class. Alluvium-fluvial has no positive cell, so it is never retained and is pooled with the reference.

\begin{table}[ht]
\centering
\small
\caption{Upper panel: agreement between the stored lithology class and four candidate assignment rules (1,899 cells; the majority rule is undefined for 10 cells with no valid raster pixel). Lower panel: lithology indicators retained in each training partition.}
\label{tab:s2}
\begin{tabular}{p{0.62\textwidth}rrr}
\toprule
Candidate rule & Cells & Matches & Agreement \\
\midrule
Raster class at the cell centroid & 1,899 & 798 & 42.0\% \\
Majority raster class within the cell & 1,889 & 858 & 45.4\% \\
Class of the contact line nearest the centroid & 1,899 & 942 & 49.6\% \\
Class of any contact line at the minimum polygon-to-contact distance & 1,899 & 1,812 & 95.4\% \\
\bottomrule
\end{tabular}


\vspace{1em}
\resizebox{\textwidth}{!}{\begin{tabular}{llll}
\toprule
Fold & Reference class & Retained indicators & Pooled with reference \\
\midrule
1 & 4 Metasedimentary & 6 Silicic; 7 Metaigneous; 16 Alluvium-Saline; 20 Water & 14 Alluvium-Fluvial \\
2 & 4 Metasedimentary & 6 Silicic; 7 Metaigneous; 16 Alluvium-Saline; 20 Water & 14 Alluvium-Fluvial \\
3 & 4 Metasedimentary & 6 Silicic; 7 Metaigneous; 16 Alluvium-Saline; 20 Water & 14 Alluvium-Fluvial \\
4 & 4 Metasedimentary & 6 Silicic; 7 Metaigneous; 16 Alluvium-Saline; 20 Water & 14 Alluvium-Fluvial \\
\bottomrule
\end{tabular}
}
\end{table}

%==========================================================================
\section{Domain distance coefficients}
\label{sec:s3}

Table~\ref{tab:s3} gives the posterior medians and 95\% equal-tailed credible intervals of the domain-varying distance coefficients for every domain and fold fit. Coefficients are per training-fold standard deviation of the predictor. CRZ\_4b, NKB, MMSB, and SRB contain no positive cells, so their coefficients are informed by background cells and partial pooling and are not estimated geological effects. MMSB is absent from the training data of the Fold~4 fit, and SRB from that of the Fold~2 fit.

\begin{table}[p]
\centering
\scriptsize
\caption{Posterior medians and 95\% credible intervals of the domain distance coefficients by fold fit. ``Contact'' is the centroid-to-internal-contact distance.}
\label{tab:s3}
\resizebox{\textwidth}{!}{\begin{tabular}{lllllr}
\toprule
Predictor & Domain & Fold & $\beta_{\mathrm{lin}}$ median [95\% CI] & $\beta_{\mathrm{sq}}$ median [95\% CI] & $P(\beta_{\mathrm{sq}}>0)$ \\
\midrule
Fault & NRB\_3a & 1 & $-1.26$ [$-1.72$, $-0.85$] & $+0.37$ [$+0.03$, $+0.72$] & 0.984 \\
Fault & NRB\_3a & 2 & $-2.59$ [$-4.41$, $-0.65$] & $+0.02$ [$-1.25$, $+2.23$] & 0.514 \\
Fault & NRB\_3a & 3 & $-1.38$ [$-2.34$, $-0.47$] & $+0.69$ [$-0.20$, $+2.04$] & 0.923 \\
Fault & NRB\_3a & 4 & $-0.68$ [$-1.74$, $+0.14$] & $+1.61$ [$+0.33$, $+2.73$] & 0.997 \\
Fault & NRB\_3b & 1 & $+0.10$ [$-0.73$, $+0.99$] & $-0.16$ [$-0.67$, $+0.23$] & 0.229 \\
Fault & NRB\_3b & 2 & $+1.29$ [$-0.92$, $+3.72$] & $-0.21$ [$-0.99$, $+0.50$] & 0.293 \\
Fault & NRB\_3b & 3 & $+0.03$ [$-0.64$, $+0.73$] & $+0.14$ [$-0.21$, $+0.48$] & 0.793 \\
Fault & NRB\_3b & 4 & $+0.12$ [$-0.48$, $+0.75$] & $+0.12$ [$-0.21$, $+0.43$] & 0.767 \\
Fault & CRZ\_4b & 1 & $-0.75$ [$-3.27$, $+1.60$] & $-0.36$ [$-2.65$, $+0.42$] & 0.232 \\
Fault & CRZ\_4b & 2 & $-1.42$ [$-5.00$, $+1.29$] & $-0.49$ [$-3.03$, $+0.55$] & 0.201 \\
Fault & CRZ\_4b & 3 & $-1.15$ [$-3.96$, $+0.77$] & $-0.28$ [$-3.36$, $+0.66$] & 0.334 \\
Fault & CRZ\_4b & 4 & $-0.54$ [$-3.21$, $+0.98$] & $-0.70$ [$-3.69$, $+0.66$] & 0.202 \\
Fault & NKB & 1 & $-0.72$ [$-3.00$, $+1.32$] & $-0.32$ [$-2.62$, $+0.56$] & 0.273 \\
Fault & NKB & 2 & $-0.92$ [$-4.55$, $+1.99$] & $-0.43$ [$-2.77$, $+0.68$] & 0.230 \\
Fault & NKB & 3 & $-0.80$ [$-3.17$, $+1.26$] & $-0.18$ [$-2.95$, $+0.74$] & 0.379 \\
Fault & NKB & 4 & $-0.42$ [$-2.83$, $+1.28$] & $-0.50$ [$-3.73$, $+0.70$] & 0.273 \\
Fault & MMSB & 1 & $-0.16$ [$-2.01$, $+2.56$] & $-0.25$ [$-2.95$, $+1.03$] & 0.359 \\
Fault & MMSB & 2 & $-0.03$ [$-2.85$, $+3.77$] & $-0.32$ [$-2.86$, $+1.37$] & 0.309 \\
Fault & MMSB & 3 & $-0.35$ [$-2.35$, $+2.48$] & $-0.06$ [$-2.95$, $+1.44$] & 0.457 \\
Fault & MMSB & 4 & $-0.32$ [$-2.56$, $+1.60$] & $+0.12$ [$-3.14$, $+3.06$] & 0.546 \\
Fault & SRB & 1 & $-0.66$ [$-3.01$, $+1.52$] & $-0.32$ [$-2.55$, $+0.49$] & 0.279 \\
Fault & SRB & 2 & $-0.49$ [$-4.49$, $+3.44$] & $-0.27$ [$-2.47$, $+1.81$] & 0.339 \\
Fault & SRB & 3 & $-0.67$ [$-3.33$, $+1.61$] & $+0.00$ [$-2.89$, $+1.59$] & 0.502 \\
Fault & SRB & 4 & $-0.36$ [$-2.80$, $+1.44$] & $-0.05$ [$-3.46$, $+2.16$] & 0.477 \\
\midrule
Contact & NRB\_3a & 1 & $-0.70$ [$-1.00$, $-0.42$] & $+0.37$ [$+0.22$, $+0.54$] & 1.000 \\
Contact & NRB\_3a & 2 & $-0.57$ [$-0.86$, $-0.29$] & $+0.35$ [$+0.17$, $+0.51$] & 1.000 \\
Contact & NRB\_3a & 3 & $-0.93$ [$-1.29$, $-0.59$] & $+0.39$ [$+0.15$, $+0.62$] & 0.999 \\
Contact & NRB\_3a & 4 & $-0.61$ [$-1.08$, $-0.11$] & $+0.38$ [$+0.14$, $+0.60$] & 0.998 \\
Contact & NRB\_3b & 1 & $-1.33$ [$-2.93$, $-0.51$] & $+0.43$ [$-0.47$, $+1.49$] & 0.885 \\
Contact & NRB\_3b & 2 & $-1.05$ [$-2.70$, $-0.29$] & $+0.33$ [$-0.70$, $+1.12$] & 0.810 \\
Contact & NRB\_3b & 3 & $-1.88$ [$-3.36$, $-1.10$] & $+0.29$ [$-0.71$, $+0.76$] & 0.787 \\
Contact & NRB\_3b & 4 & $-1.90$ [$-3.48$, $-0.96$] & $+0.43$ [$-0.52$, $+1.21$] & 0.878 \\
Contact & CRZ\_4b & 1 & $-0.70$ [$-2.66$, $+2.05$] & $+0.34$ [$-1.71$, $+1.81$] & 0.739 \\
Contact & CRZ\_4b & 2 & $-0.52$ [$-2.18$, $+2.00$] & $+0.26$ [$-1.89$, $+1.40$] & 0.696 \\
Contact & CRZ\_4b & 3 & $-0.80$ [$-2.48$, $+2.31$] & $+0.22$ [$-2.00$, $+1.13$] & 0.664 \\
Contact & CRZ\_4b & 4 & $-0.54$ [$-2.49$, $+2.52$] & $+0.31$ [$-1.88$, $+1.36$] & 0.719 \\
Contact & NKB & 1 & $-0.78$ [$-2.97$, $+1.42$] & $+0.32$ [$-1.64$, $+1.47$] & 0.736 \\
Contact & NKB & 2 & $-0.60$ [$-2.61$, $+1.90$] & $+0.30$ [$-1.55$, $+1.76$] & 0.724 \\
Contact & NKB & 3 & $-1.15$ [$-3.53$, $+1.00$] & $+0.21$ [$-1.97$, $+0.90$] & 0.660 \\
Contact & NKB & 4 & $-0.90$ [$-3.34$, $+1.41$] & $+0.31$ [$-1.69$, $+1.32$] & 0.734 \\
Contact & MMSB & 1 & $-0.55$ [$-2.19$, $+2.39$] & $+0.28$ [$-1.96$, $+1.46$] & 0.687 \\
Contact & MMSB & 2 & $-0.51$ [$-2.02$, $+2.20$] & $+0.27$ [$-1.88$, $+1.45$] & 0.685 \\
Contact & MMSB & 3 & $-0.82$ [$-2.50$, $+2.25$] & $+0.24$ [$-2.06$, $+1.25$] & 0.679 \\
Contact & MMSB & 4 & $-0.80$ [$-3.43$, $+1.99$] & $+0.34$ [$-1.51$, $+1.66$] & 0.753 \\
Contact & SRB & 1 & $-0.72$ [$-2.78$, $+1.95$] & $+0.34$ [$-1.68$, $+1.70$] & 0.743 \\
Contact & SRB & 2 & $-0.60$ [$-2.61$, $+2.02$] & $+0.30$ [$-1.68$, $+1.92$] & 0.730 \\
Contact & SRB & 3 & $-0.98$ [$-3.03$, $+1.98$] & $+0.26$ [$-1.76$, $+1.26$] & 0.700 \\
Contact & SRB & 4 & $-0.72$ [$-3.07$, $+2.23$] & $+0.33$ [$-1.64$, $+1.73$] & 0.744 \\
\bottomrule
\end{tabular}
}
\end{table}

%==========================================================================
\section{Direction of distance associations across representations}
\label{sec:s4}

Table~\ref{tab:s4} reports the posterior probability that the linear distance coefficient is negative, $P(\beta_{\mathrm{lin}}<0)$, in NRB\_3a and NRB\_3b under each of the four distance representations. In the quadratic models (A and B), the linear coefficient is the slope of the logit surface at the training-mean transformed distance. In the linear models (C and D), it is the slope everywhere. A value near 1 means that probability decreases with distance from the feature.

\begin{table}[ht]
\centering
\small
\caption{Posterior probability of a negative linear distance coefficient by domain, predictor, representation, and fold fit.}
\label{tab:s4}
\begin{tabular}{lllrrrr}
\toprule
Domain & Predictor & Model & Fold 1 & Fold 2 & Fold 3 & Fold 4 \\
\midrule
NRB\_3a & Contact & A (raw, quadratic) & 1.000 & 1.000 & 1.000 & 0.993 \\
NRB\_3a & Contact & B (log, quadratic) & 0.998 & 0.980 & 1.000 & 0.885 \\
NRB\_3a & Contact & C (raw, linear) & 0.996 & 0.969 & 1.000 & 0.584 \\
NRB\_3a & Contact & D (log, linear) & 1.000 & 1.000 & 1.000 & 0.962 \\
\addlinespace
NRB\_3a & Fault & A (raw, quadratic) & 1.000 & 0.997 & 0.999 & 0.942 \\
NRB\_3a & Fault & B (log, quadratic) & 1.000 & 1.000 & 1.000 & 0.992 \\
NRB\_3a & Fault & C (raw, linear) & 1.000 & 1.000 & 1.000 & 1.000 \\
NRB\_3a & Fault & D (log, linear) & 1.000 & 1.000 & 1.000 & 1.000 \\
\addlinespace
NRB\_3b & Contact & A (raw, quadratic) & 1.000 & 0.995 & 1.000 & 1.000 \\
NRB\_3b & Contact & B (log, quadratic) & 0.999 & 0.979 & 1.000 & 1.000 \\
NRB\_3b & Contact & C (raw, linear) & 1.000 & 0.985 & 1.000 & 1.000 \\
NRB\_3b & Contact & D (log, linear) & 1.000 & 0.999 & 1.000 & 1.000 \\
\addlinespace
NRB\_3b & Fault & A (raw, quadratic) & 0.405 & 0.121 & 0.459 & 0.356 \\
NRB\_3b & Fault & B (log, quadratic) & 0.539 & 0.639 & 0.262 & 0.230 \\
NRB\_3b & Fault & C (raw, linear) & 0.738 & 0.015 & 0.108 & 0.080 \\
NRB\_3b & Fault & D (log, linear) & 0.727 & 0.180 & 0.333 & 0.300 \\
\bottomrule
\end{tabular}

\end{table}

%==========================================================================
\section{Probability accuracy and calibration}
\label{sec:s5}

Table~\ref{tab:s5} compares the DVHM Brier score with two prevalence baselines and reports logistic calibration diagnostics. The out-of-fold baseline gives each held-out cell the positive rate of its own training partition and is the fair comparator. The pooled-prevalence baseline uses the label frequency of the held-out cells and is shown only for reference. The calibration slope is the coefficient of a logistic regression of the label on $\mathrm{logit}(\bar p_i)$. A slope below 1 means that predictions are too extreme. Calibration-in-the-large is the intercept when $\mathrm{logit}(\bar p_i)$ enters as an offset.

\begin{table}[ht]
\centering
\small
\caption{Upper panel: Brier skill and calibration of pooled out-of-fold DVHM probabilities (1,872 cells, 138 positive). Lower panel: reliability by decile of predicted probability.}
\label{tab:s5}
\begin{tabular}{lr}
\toprule
Quantity & Value \\
\midrule
DVHM Brier score & 0.0688 \\
Baseline: training-partition prevalence (out-of-fold) & 0.0695 \\
Baseline: pooled prevalence & 0.0683 \\
Brier skill vs.\ training-partition prevalence & $+0.010$ \\
Brier skill vs.\ pooled prevalence & $-0.007$ \\
Calibration slope (logistic, on logit $\bar p$) & 0.605 \\
Calibration intercept (with slope) & $-0.836$ \\
Calibration-in-the-large (logit offset) & $+0.049$ \\
Mean predicted probability / observed prevalence & 0.0708 / 0.0737 \\
\midrule
\multicolumn{2}{l}{\emph{By held-out fold: Brier (baseline) and skill}} \\
Fold 1 & 0.0185 (0.0241); $+0.233$ \\
Fold 2 & 0.0627 (0.0620); $-0.012$ \\
Fold 3 & 0.0765 (0.0729); $-0.050$ \\
Fold 4 & 0.1175 (0.1191); $+0.013$ \\
\bottomrule
\end{tabular}


\vspace{1em}
\begin{tabular}{rrrrr}
\toprule
Decile & Cells & Mean predicted & Observed rate & Positive cells \\
\midrule
1 & 188 & 0.005 & 0.000 & 0 \\
2 & 187 & 0.010 & 0.011 & 2 \\
3 & 187 & 0.016 & 0.021 & 4 \\
4 & 187 & 0.022 & 0.043 & 8 \\
5 & 187 & 0.033 & 0.053 & 10 \\
6 & 187 & 0.050 & 0.102 & 19 \\
7 & 187 & 0.069 & 0.107 & 20 \\
8 & 187 & 0.093 & 0.139 & 26 \\
9 & 187 & 0.133 & 0.118 & 22 \\
10 & 188 & 0.278 & 0.144 & 27 \\
\bottomrule
\end{tabular}

\end{table}

%==========================================================================
\section{Posterior uncertainty by domain and training support}
\label{sec:s6}

Table~\ref{tab:s6} summarizes the width of the cell-level 95\% credible interval of the conditional probability by domain and by the training support of the cell's domain in its fold fit.

\begin{table}[ht]
\centering
\small
\caption{Width of the 95\% credible interval of the conditional probability for the 1,872 held-out cells. $^{\dagger}$Domain without positive cells. ``Domain absent from training'' covers the 16 MMSB cells (Fold~4) and the SRB cell (Fold~2).}
\label{tab:s6}
\begin{tabular}{lrrrrrr}
\toprule
Group & Cells & Median & 90th pct. & Max & $>0.2$ & $>0.6$ \\
\midrule
NRB\_3a & 1127 & 0.098 & 0.281 & 0.975 & 18.2\% & 2.9\% \\
NRB\_3b & 550 & 0.084 & 0.332 & 0.937 & 17.5\% & 6.0\% \\
CRZ\_4b$^{\dagger}$ & 154 & 0.102 & 0.890 & 0.991 & 37.0\% & 18.2\% \\
NKB$^{\dagger}$ & 24 & 0.382 & 1.000 & 1.000 & 66.7\% & 37.5\% \\
MMSB$^{\dagger}$ & 16 & 0.656 & 0.800 & 0.869 & 100.0\% & 75.0\% \\
SRB$^{\dagger}$ & 1 & 0.977 & 0.977 & 0.977 & 100.0\% & 100.0\% \\
\midrule
Domain absent from training & 17 & 0.656 & 0.851 & 0.977 & 100.0\% & 76.5\% \\
In training, no positive cells & 178 & 0.109 & 0.930 & 1.000 & 41.0\% & 20.8\% \\
In training, with positive cells & 1677 & 0.095 & 0.284 & 0.975 & 17.9\% & 3.9\% \\
\bottomrule
\end{tabular}

\end{table}

%==========================================================================
\section{Hierarchical population and scale parameters}
\label{sec:s7}

Table~\ref{tab:s7} summarizes the population means and between-domain scales of the V14 DVHM fold fits, derived directly from the posterior samples. These are not the hyperparameters of the separate full-data descriptive model used for exploratory heterogeneity analysis. With only two domains contributing positive cells, the scales are weakly identified, and their upper credible limits stay close to the HalfNormal(1) prior.

\begin{table}[ht]
\centering
\small
\caption{Posterior medians and 95\% credible intervals of the population means ($\mu$), between-domain standard deviations ($\sigma$), and global gravity coefficient of the DVHM, by fold fit.}
\label{tab:s7}
\resizebox{\textwidth}{!}{\begin{tabular}{lllll}
\toprule
Parameter & Fold 1 & Fold 2 & Fold 3 & Fold 4 \\
\midrule
$\mu_\alpha$ & $-3.52$ [$-4.51$, $-2.51$] & $-4.22$ [$-5.34$, $-2.78$] & $-4.27$ [$-5.21$, $-3.20$] & $-4.58$ [$-5.50$, $-3.37$] \\
$\sigma_\alpha$ & $+0.57$ [$+0.02$, $+1.93$] & $+0.85$ [$+0.05$, $+2.32$] & $+0.60$ [$+0.04$, $+1.85$] & $+0.56$ [$+0.03$, $+1.89$] \\
$\mu_{f,\mathrm{lin}}$ & $-0.52$ [$-1.58$, $+0.66$] & $-0.53$ [$-1.93$, $+1.12$] & $-0.63$ [$-1.73$, $+0.56$] & $-0.34$ [$-1.48$, $+0.62$] \\
$\sigma_{f,\mathrm{lin}}$ & $+0.89$ [$+0.29$, $+2.07$] & $+1.59$ [$+0.42$, $+2.89$] & $+0.93$ [$+0.22$, $+2.15$] & $+0.69$ [$+0.06$, $+1.89$] \\
$\mu_{f,\mathrm{sq}}$ & $-0.16$ [$-1.25$, $+0.41$] & $-0.27$ [$-1.36$, $+0.65$] & $+0.01$ [$-1.29$, $+0.76$] & $+0.06$ [$-1.31$, $+1.03$] \\
$\sigma_{f,\mathrm{sq}}$ & $+0.57$ [$+0.08$, $+1.76$] & $+0.53$ [$+0.02$, $+2.05$] & $+0.73$ [$+0.03$, $+2.07$] & $+1.16$ [$+0.26$, $+2.44$] \\
$\mu_{l,\mathrm{lin}}$ & $-0.72$ [$-1.66$, $+0.50$] & $-0.58$ [$-1.48$, $+0.65$] & $-0.96$ [$-1.89$, $+0.41$] & $-0.78$ [$-1.84$, $+0.58$] \\
$\sigma_{l,\mathrm{lin}}$ & $+0.72$ [$+0.05$, $+2.02$] & $+0.64$ [$+0.04$, $+2.00$] & $+0.87$ [$+0.13$, $+2.18$] & $+0.97$ [$+0.23$, $+2.31$] \\
$\mu_{l,\mathrm{sq}}$ & $+0.33$ [$-0.82$, $+1.02$] & $+0.28$ [$-0.84$, $+1.02$] & $+0.25$ [$-1.03$, $+0.73$] & $+0.33$ [$-0.77$, $+0.99$] \\
$\sigma_{l,\mathrm{sq}}$ & $+0.41$ [$+0.02$, $+1.67$] & $+0.39$ [$+0.02$, $+1.63$] & $+0.34$ [$+0.01$, $+1.60$] & $+0.37$ [$+0.02$, $+1.58$] \\
$\beta_g$ & $-0.29$ [$-0.55$, $-0.04$] & $-0.13$ [$-0.37$, $+0.09$] & $-0.46$ [$-0.74$, $-0.20$] & $-0.65$ [$-1.00$, $-0.30$] \\
\midrule
\multicolumn{5}{l}{Minimum bulk ESS 1,219; maximum $\hat R$ 1.000 across these parameters and folds.} \\
\bottomrule
\end{tabular}
}
\end{table}

%==========================================================================
\section{Stationary-point robustness in the informative domains}
\label{sec:s8}

Table~\ref{tab:s8} lists the 16 representation-robustness cases in NRB\_3a and NRB\_3b, the two domains with positive cells. $P_A$ and $P_B$ are the posterior probabilities of positive curvature under Models~A and~B. $D^*_A$ and $D^*_B$ are posterior median stationary points in kilometres. C1--C3 are the three criteria of the main text.

\begin{table}[ht]
\centering
\small
\caption{Representation-robustness criteria for the 16 informative cases.}
\label{tab:s8}
\resizebox{\textwidth}{!}{\begin{tabular}{lllrrcrrccc}
\toprule
Fold & Domain & Predictor & $P_A$ & $P_B$ & C1 & $D^*_A$ & $D^*_B$ & C2 & C3 & All \\
\midrule
1 & NRB\_3a & Fault & 0.98 & 0.11 & -- & 37.1 & -0.4 & -- & -- & -- \\
2 & NRB\_3a & Fault & 0.51 & 0.11 & -- & 24.2 & 0.1 & -- & \checkmark & -- \\
3 & NRB\_3a & Fault & 0.92 & 0.08 & -- & 48.0 & 0.0 & -- & \checkmark & -- \\
4 & NRB\_3a & Fault & 1.00 & 0.76 & -- & 34.3 & 97.0 & -- & -- & -- \\
1 & NRB\_3b & Fault & 0.23 & 0.08 & -- & 22.1 & 7.0 & -- & \checkmark & -- \\
2 & NRB\_3b & Fault & 0.29 & 0.89 & -- & 77.9 & 14.1 & -- & \checkmark & -- \\
3 & NRB\_3b & Fault & 0.79 & 0.55 & -- & 34.9 & 12.4 & -- & \checkmark & -- \\
4 & NRB\_3b & Fault & 0.77 & 0.70 & \checkmark & 33.1 & 10.3 & -- & \checkmark & -- \\
1 & NRB\_3a & Contact & 1.00 & 0.83 & -- & 31.1 & 35.8 & \checkmark & \checkmark & -- \\
2 & NRB\_3a & Contact & 1.00 & 0.83 & -- & 30.5 & 28.5 & \checkmark & \checkmark & -- \\
3 & NRB\_3a & Contact & 1.00 & 0.56 & -- & 28.0 & 28.2 & \checkmark & \checkmark & -- \\
4 & NRB\_3a & Contact & 1.00 & 0.89 & \checkmark & 29.6 & 16.8 & -- & \checkmark & -- \\
1 & NRB\_3b & Contact & 0.89 & 0.37 & -- & 33.8 & 0.5 & -- & \checkmark & -- \\
2 & NRB\_3b & Contact & 0.81 & 0.74 & \checkmark & 32.4 & 27.7 & \checkmark & \checkmark & \checkmark \\
3 & NRB\_3b & Contact & 0.79 & 0.15 & -- & 38.0 & -0.2 & -- & -- & -- \\
4 & NRB\_3b & Contact & 0.88 & 0.23 & -- & 42.5 & -0.0 & -- & -- & -- \\
\midrule
\multicolumn{11}{l}{All 48 cases: C1 15, C2 4, C3 18, all three 1. Informative 16: C1 3, C2 4, C3 12, all three 1.} \\
\bottomrule
\end{tabular}
}
\end{table}

\end{document}
